\chapter{Exam-to-Textbook Bridge}
\label{app:exam-bridge}

\section{From the Midterm}

\begin{itemize}
\item Differential forms and vector calculus become the language of \(1\)-forms
  and residues in \cref{ch:exam-background}.
\item The winding form is the prototype for a closed local object that is not
  globally exact.
\item Cauchy--Green and Stokes explain why deleting small disks around
  singularities leaves residue contributions.
\item The Euler characteristic computation for the torus supplies \(g=1\), used
  in \cref{ch:riemann-roch}.
\end{itemize}

\section{From the Final}

\begin{itemize}
\item Compact Riemann surfaces and meromorphic maps are the setting of
  \cref{ch:riemann-surfaces}.
\item The theorem that principal divisors have degree zero appears as
  \cref{prop:principal-degree-zero}.
\item Riemann--Roch supplies the dimension count \(\ell(nO)=n\) in
  \cref{ch:riemann-roch}.
\item The final problem on the Weierstrass equation becomes
  \cref{thm:wp-ode-new}.
\item Base point freeness and separation are the background for why a
  three-dimensional linear system can embed a genus-one curve in projective
  space.
\end{itemize}

\begin{keyidea}{The exam arc}
The midterm gives the obstruction philosophy; the final gives the compact
Riemann surface technology.  Together they lead to the equivalence
\(\C/\Lambda\cong\Pic^0(E)\cong E(\C)\) as complex Lie groups, not merely as
abstract abelian groups.
\end{keyidea}
